\documentclass[../Main.tex]{subfiles}

\begin{document}
\section{Defining a Field}
A field is a set that permits addition and multiplication, including inverses (except the multiplicative inverse of $0$), and has the distribution property:
\begin{equation*}
    a(b+c) = ab + ac.
\end{equation*}
\section{Defining a Vector Space}
\begin{definition}{Vector space}
    A \underline{vector space} over a field $\F$ is an Abelian group $(V, +, \zv)$ equipped with scalar multiplication:
    \begin{align*}
        \F \times V &\mapsto V\\
        (\lambda, v) &\mapsto \lambda v
    \end{align*}
    that satisfies, for $\lambda, \mu \in \F$ and $v, w \in V$:
    \begin{enumerate}
        \item $(\lambda + \mu) v = \lambda v + \mu v$
        \item $\lambda(v + w) = \lambda v + \lambda w$
        \item $\lambda(\mu v) = (\lambda \times \mu) v$
        \item $1 v = v$.
    \end{enumerate}
\end{definition}
\section{Subspaces}
\begin{definition}{Subspace}
    A \underline{subspace} $U$ of a vector space $V$ is a subset of $V$ which is a vector space over $\F$ with the same operations as $V$.
\end{definition}
\begin{proposition}[Subspace test]
    Let $V$ be a vector space over $\F$. Let $U \subseteq V$. Then $U$ is a subspace of $V$ if and only if:
    \begin{enumerate}
        \item $U \neq \emptyset$,
        \item $\forall u, v \in U,~\forall \lambda \in \F,~u + \lambda w \in U$.
    \end{enumerate}
    \label{propSubspaceTest}
\end{proposition}
\begin{lemma}
    For vector spaces $U, W \leq V$, the intersection $U \cap W$ is a subspace of $V$.
    \label{lemIntersecIsSubspace}
\end{lemma}
\begin{proof}
    Apply \thmref{propSubspaceTest}.
\end{proof}
\begin{definition}{Subspace sum}
    The \underline{subspace sum} of $U, W \leq V$ is:
    \begin{equation*}
        U + W = \subsetselect{u + w}{u \in U, w \in W}
    \end{equation*}
\end{definition}
\begin{proof}
    Apply \thmref{propSubspaceTest}.
\end{proof}
\section{Linear Maps}
\begin{definition}{Linear map}
    A \underline{linear map} from $V$ to $W$ is a function $\alpha : V \mapsto W$ such that for all $u, v \in V$ and $\lambda \in \F$, $\alpha(u + \lambda v) = \alpha(u) + \lambda \alpha(v)$.
\end{definition}
\begin{definition}{Isomorphism}
    An \underline{isomorphism} is a linear map that is bijective. If such a linear map exists for vector spaces $V$ and $W$, they are then \underline{isomorphic} and we write $V \isom W$.
\end{definition}
\begin{proposition}
    Let $\alpha : V \mapsto W$ be an isomorphism. Then so is $\alpha^{-1}$.
    \label{propIsomInvertible}
\end{proposition}
\begin{proof}
    We know that $\alpha^{-1}$ is a bijection so it remains to show linearity. To do this, let $w_1 = \alpha(v_1)$ and the same for $w_2$ and then apply the definition.
\end{proof}
\begin{definition}{Image}
    The \underline{image} of a linear map $\alpha : V \mapsto W$ is the set:
    \begin{equation*}
        \im(\alpha) = \subsetselect{\alpha(v) \in W}{v \in V}
    \end{equation*}
\end{definition}
\begin{definition}{kernel}
    The \underline{kernel} of a linear map $\alpha : V \mapsto W$ is the set:
    \begin{equation*}
        \ker(\alpha) = \subsetselect{v \in V}{\alpha(v) = \zv}
    \end{equation*}
\end{definition}
\begin{propositions}{
        Let $\alpha : V \mapsto W$ be a linear map.
        \label{propsImageKernel}
    }
    \item $\ker(\alpha) \leq V$
    \item $\im(\alpha) \leq W$
    \item $\alpha$ is surjective if and only if $\im(\alpha) = W$
    \item $\alpha$ is injective if and only if $\ker(\alpha) = \{\zv\}$.
\end{propositions}
\begin{proof}
    The first two are an application of \thmref{propSubspaceTest}. Then the third is very simple from the definition of surjectivity, the fourth is by considering the difference of two elements of the kernel.
\end{proof}
\section{Quotient Spaces}
\begin{definition}{Left coset}
    Let $V$ be a vector space over $\F$ with a subspace $W$. The \underline{left coset} of $W$ by an element of $V$ is:
    \begin{equation*}
        v + W = \subsetselect{v + w}{w \in V}
    \end{equation*}
\end{definition}
The set of left cosets is $V / W$.
\begin{definition}{Quotient space}
    The \underline{quotient space} of a vector space $V$ by a subspace $W$ is the vector space $V / W$ with the operations:
    \begin{enumerate}
        \item $(u + W) + (v + W) = (u + v) + W$
        \item $\lambda (v + W) = (\lambda v) + W$.
    \end{enumerate}
\end{definition}
\begin{proposition}
    The quotient space is a well-defined vector space.
    \label{propQuotientSpace}
\end{proposition}
\begin{proof}
    The proof is similar to that of IA Groups.
\end{proof}
\begin{proposition}[Quotient Map]
    There exists a well-defined function:
    \begin{align*}
        \pi_W : V &\mapsto V / W\\
        v &\mapsto v + W
    \end{align*}
    which is surjective, and has kernel $W$.
    \label{propQuotientMap}
\end{proposition}
\begin{proof}
    This is again a simple proof and the statements can be tackled one at a time.
\end{proof}
\begin{theorem}[First Isomorphism Theorem]
    Let $\alpha : V \mapsto W$ be a linear map. There exists an isomorphism:
    \begin{align*}
        \overline{\alpha} : V / \ker(\alpha) &\mapsto \im(\alpha)\\
        v + \ker(\alpha) &\mapsto \alpha(v)
    \end{align*}
    \label{thmIsomorphism}
\end{theorem}
\begin{proof}
    We can show well-defined and injective at once by chaining equivalences to show $u + \ker(\alpha) = v + \ker(\alpha) \iff \overline{\alpha}(u + \ker(\alpha)) = \overline{\alpha}(v + \ker(\alpha))$. It's then easy to show that $\overline{\alpha}$ is surjective, and that it is linear.
\end{proof}
\end{document}